\documentclass[11pt]{article}
\usepackage{amsmath,amssymb,amsthm}
\newtheorem{lemma}{Lemma}
\newtheorem{claim}{Claim}
\theoremstyle{definition}\newtheorem{definition}{Definition}
\theoremstyle{remark}\newtheorem*{remark}{Remark}
\title{Idea A made precise: line passes on the fine field of one axis}
\date{Draft, October 8, 2026}
\begin{document}
\maketitle

This replaces the paragraph ``Idea A'' of \texttt{notes/stack-notes.tex}. It refers to the
pinned upstream sections \texttt{02-streams.tex} (machine model,
Lemma~elementary-streams and its counter paragraph) and \texttt{07-resampling.tex}
(lem:tensor-resampling), to \texttt{notes/segmented-inverse.tex} and to PR~\#5's
chirped-correlation blocks. A machine-level simulation of every data movement below is
\texttt{scripts/check\_fine\_field.py}; it checks outputs against the linewise maps
and checks the movement bounds of Lemma~\ref{lem:fine} case by case.

\section{Setting}

Fix an axis $i$ of the padded box. Its box length is $n=t_i\in\{r/2,r\}$, a power of two.
Collapse the address fields before axis $i$ to one field $X$ and those after it to one
field $Y$ (upstream's collapsing convention). An entry is a record of $b=O(p)$ bits.
Movement is counted in records; a record move costs $b$ head steps. Put
$V=|X|\,n\,|Y|$ (records), choose $f$ with $2^f\mid n$, write the axis coordinate as
\[ j=J2^f+u,\qquad 0\le u<2^f,\quad 0\le J<n_s:=n/2^f, \]
and call $(X,J,Y,u)$, lexicographic with $X$ most significant, the \emph{slab order}. For
fixed $X$ and $J$ the $|Y|2^f$ records with that $(X,J)$ form the \emph{slab} $(X,J)$; for
fixed $Y$ in it, the $2^f$ consecutive records form the \emph{slab row} $(X,J,Y)$. Put
$B=|Y|2^f$ and $G=n_sB$ (records per $X$-group).

A line $(X,Y)$ is \emph{valid} when every coordinate held in $X$ and $Y$ lies below its
current bound ($s_h$ or $t_h$, per the shape descriptor). Invalid lines are zero and must be
mapped to zero lines (upstream proof of lem:tensor-resampling).

\begin{definition}[Slab-local sweep]\label{def:sweep}
A \emph{forward sweep} on a line of box length $n$ is given by
\begin{itemize}
\item an input period $\pi\in\{s_i,t_i\}$: the input $x$ is read only through the virtual
line $\tilde x_j=x_{j\bmod\pi}$, $j\in\mathbb Z$, so box entries $j\in[\pi,n)$ are never read;
\item an output count $m\le n$: outputs $k\in[0,m)$ are computed and $k\in[m,n)$ are written $0$;
\item a non-decreasing centre $c:\mathbb Z\to\mathbb Z$ with $c(0)=0$ and
$c(k+a)-c(k)\le\sigma a+1$ for $a\ge0$, where $\sigma\le5/4$, and a half-width $W$ such that
output $k$ reads $\tilde x$ only on $[c(k)-W,c(k)+W]$;
\item a state of $\Sigma$ records per line, an auxiliary row of $\xi$ records per slab row,
and per-slab data $D_J$ that depend on $J$ and the axis only (not on $X$, $Y$);
\item a kernel $K_J$ that, given $\tilde x$ on the window of Claim~\ref{cl:window}, the state
after slab $J-1$ and $D_J$, returns the $2^f$ outputs of slab $J$, the state after slab $J$ and
the auxiliary row of slab $J$.
\end{itemize}
A \emph{backward sweep} is the same with $J$ decreasing, its state seeded from the final
forward state, and its kernel reading the auxiliary row that the forward sweep wrote for the
same slab row. The \emph{linewise map} is the sweep run on one line in natural order.
\end{definition}

The passes Idea A must carry are instances (Section~\ref{sec:instances}): $\mathsf S'$ and
$A$ ($\pi=s_i$, $m=t_i$, $c(k)=\lfloor s_ik/t_i\rfloor$), $B_0$ ($\pi=t_i$, $m=s_i$,
$c(k)=\lfloor t_ik/s_i\rfloor$, so $\sigma=t_i/s_i\le5/4$), $N$, $M^{-1}$ and $E$
($\pi=s_i$, $m=s_i$, $c(k)=k$), and the Woodbury forward and backward sweeps of the segmented
inverse.

\section{The lemma}

\begin{lemma}[Fine-field line passes]\label{lem:fine}
Let a forward sweep satisfy
\begin{enumerate}
\item[(W1)] $2^f\ge4W+4$;
\item[(W2)] $\pi\ge3\cdot2^f$;
\item[(W3)] $\Sigma\le2^f$ and $\xi\le2^f$.
\end{enumerate}
On a fixed set of work tapes (independent of $d$, $|X|$, $|Y|$, $n$) the following hold, with
all counters, rewinds and returns included.
\begin{enumerate}
\item[(a)] \emph{Exposure.} $(X,J,u,Y)\to(X,J,Y,u)$ and its inverse are $f$ digit moves each and
cost exactly $6fV$ record moves each.
\item[(b)] \emph{Forward pass.} The schedule of Section~\ref{sec:fwd} writes, in slab order,
exactly the outputs of the linewise map on every valid line and zeros elsewhere, and moves at most
\[
 28V+11|X|(n_e-n_s)B+\frac{4\Sigma+2\xi+2}{2^f}V+4|X||Y|\Sigma
\]
records, where $n_e\le n_s+4$ is the number of extended slabs.
\item[(c)] \emph{Backward pass.} The schedule of Section~\ref{sec:bwd} is likewise exact and moves
at most $11V+\frac{3\xi+6\Sigma+2}{2^f}V+4|X||Y|\Sigma$ records.
\item[(d)] \emph{Setup.} The per-slab data are computed once per axis and read in lockstep with
$J$; the validity bitmap costs at most $2|Y|$ moves to build per pass. Kernel work is the
linewise work plus at most two partial blocks per slab row.
\end{enumerate}
Hence an axis with $P_f$ forward-type and $P_b$ backward-type passes costs
\[
 b\,V\bigl(12f+28P_f+11P_b+o(1)\bigr)\ \text{head steps}
\]
plus the line algorithm's own arithmetic, when $n_s\to\infty$; for every axis this is
$O(Vb(f+P))$ with $P=P_f+P_b=O(1)$.
\end{lemma}

The constants are worst-case bounds; the simulation's measured values are listed at the end.

\section{Proof}

\subsection{Exposure and restoration}
Write $u=(u_{f-1}\dots u_0)$ in binary. Before step $\iota=0,\dots,f-1$ the order is
$(X,J,u_{f-1}\dots u_\iota,Y,u_{\iota-1}\dots u_0)$. Step $\iota$ moves digit $u_\iota$ past $Y$:
with prefix $P=|X|n_s2^{f-1-\iota}$ and suffix $S=2^\iota$ it maps
$(P,u_\iota,Y,S)\to(P,Y,u_\iota,S)$, which is item~2 of Lemma~elementary-streams. Implemented as
split on $u_\iota$ into two stream tapes ($V$ reads, $V$ writes), return of both streams ($V$ in
total), merge in output order ($V$ reads, $V$ writes) and return of the output ($V$), it costs
$6V$. Restoration runs the inverse moves for $\iota=f-1,\dots,0$. When $Y$ is empty (the last
axis) nothing moves.

\subsection{The window}
\begin{claim}\label{cl:window}
Put $J'(J)=\lfloor(c(J2^f)-W)/2^f\rfloor$. Under (W1), every input read by an output of slab $J$
lies in virtual slabs $J',J'+1,J'+2$, and $0\le J'(J+1)-J'(J)\le2$.
\end{claim}
\begin{proof}
The inputs of slab $J$ lie in $[c(J2^f)-W,\,c(J2^f+2^f-1)+W]$, an interval of length at most
$\sigma(2^f-1)+2W+2\le\tfrac54 2^f+\tfrac12 2^f<2\cdot2^f+1$ by (W1). It starts at offset
$o\in[0,2^f)$ from $J'2^f$, so it ends before $J'2^f+o+2\cdot2^f+1\le(J'+3)2^f$. Monotonicity
of $c$ gives $J'(J+1)\ge J'(J)$, and $c((J+1)2^f)-c(J2^f)\le\sigma2^f+1<2\cdot2^f$ gives the step
bound.
\end{proof}

\begin{remark}[The notes' window was wrong]
\texttt{stack-notes.tex} centres the window at $J'=\lfloor J/\sigma\rfloor$ and reads slabs
$J'-1,J',J'+1$. The inputs of slab $J$ end near $(J+1)2^f/\sigma+W$, which reaches slab $J'+2$
whenever $\{J/\sigma\}>1-1/\sigma-W2^{-f}$, a positive fraction of $J$ since $\theta$ is small.
For $(s,t,2^f,W)=(1009,1024,64,8)$, 6 of 16 slabs fail (\texttt{scripts/check\_fine\_field.py}). Anchoring at the
lowest input, as above, keeps three slabs.
\end{remark}

\subsection{Virtual slabs and the wrap splice}
For an X-group, the \emph{extended array} holds virtual slabs $v=v_{\rm lo},\dots,v_{\rm hi}$ in
order $(v,Y,u)$, where entry $(v,Y,u)$ is $\tilde x_{v2^f+u}$ of line $(X,Y)$. Here
$v_{\rm lo}=J'(0)=\lfloor-W/2^f\rfloor=-1$ and $v_{\rm hi}=\max\{J'(J)+2: J2^f<m\}$. In
every instance $c(m-1)<\pi$, so $J'\le n_\pi$, $v_{\rm hi}\le n_\pi+2\le n_s+2$ and
$n_e\le n_s+4$; and $(v_{\rm hi}+1)2^f\le\pi+3\cdot2^f\le2\pi$ by (W2).

\begin{claim}\label{cl:splice}
Fix $v$ and $\rho=\pi\bmod2^f$, $n_\pi=\lfloor\pi/2^f\rfloor$. Row $(v,Y)$ is drawn from at most
two physical slab rows $(a_v,Y)$, $(b_v,Y)$, each as one run of increasing $u$:
\begin{itemize}
\item $0\le v<n_\pi$: the row is physical row $(v,Y)$ (aligned copy);
\item $v=n_\pi$, $\rho\ne0$: $u<\rho$ from row $(n_\pi,Y)$ at $u$; $u\ge\rho$ from row $(0,Y)$
at $u-\rho$ (the splice at $u=\pi\bmod2^f$);
\item $v<0$ or $v>n_\pi$: the virtual range maps to one physical interval of length $2^f$
starting at $v2^f+u\bmod\pi$, which meets two consecutive slabs, split at the shift $\rho$.
\end{itemize}
When $\rho=0$ (the input period of $B_0$ is $t_i$) every virtual row is an aligned copy.
\end{claim}
\begin{proof}
$[v2^f,(v+1)2^f)$ has length $2^f<\pi$ and lies in $[-\pi,2\pi)$, so it contains at most one
multiple of $\pi$. If none, its image mod $\pi$ is one interval of length $2^f$, which meets at
most two slabs, each in an increasing run. If it contains $\pi$ it is $v=n_\pi$, its part below
$\pi$ is $[n_\pi2^f,\pi)$ inside slab $n_\pi$, and its part above maps to $[0,2^f-\rho)$ inside
slab~$0$. It cannot contain $0$ unless $v=0$, where the range starts at $0$.
\end{proof}

\paragraph{Building it.} Tape IN holds the array; first copy the X-group to tape CP ($G$ reads,
$G$ writes, one return). For each $v$ park IN at row $(a_v,0)$ and CP at row $(b_v,0)$. For
$Y=0,\dots,|Y|-1$ read the two runs in $u$ order, each head finishing its slab row, and write
every value to E0, E1, E2 in the same step. Each head advances exactly $2^f$ per $Y$, so $B$ per
virtual slab. A single head would pay $|Y|2^f$ per row to jump between the two physical slabs,
i.e.\ $\Theta(|Y|B)$ per spliced slab; the second head is what makes the splice linear. The
sequence of $a_v$ is $n_\pi-1$ (for $v=-1$), $0,1,\dots,n_\pi$, then small indices again, so IN
seeks across the group at most four times and CP at most twice.

\subsection{Forward schedule}\label{sec:fwd}
Tapes: IN, CP, E0, E1, E2, OUT, ST (state), VB (validity bitmap over $Y$), AX (auxiliary rows),
FIN (final states), BUF (window buffer), DJ (per-slab data), descriptor and counter tapes.

\emph{Validity.} VB holds $\prod_h[y_h<\text{bound}_h]$ over $Y$ in $Y$ order. Build it from the
innermost $Y$ axis outward: the bitmap for axes $h..$ is the concatenation over digit values of
copies of the bitmap for axes $h+1..$, or of zeros above the bound. Each level costs twice its
size and sizes at least double, so building costs at most $2|Y|$ per tape. $X$-validity is
computed once per group by an odometer, $O(L)$ against a group of $G\ge n$ records.

For each $X$ in order:
\begin{enumerate}
\item If $X$ is invalid, advance IN by $G$ and write $n_sB$ zeros to OUT, $n_s|Y|\xi$ to AX and
$|Y|\Sigma$ to FIN.
\item Otherwise build the extended array and write $|Y|\Sigma$ zero states to ST.
\item For $J=0,\dots,n_s-1$: if $J2^f\ge m$ write a zero slab and zero auxiliary rows. Otherwise
seek E$r$ to row $(J'(J)+r,0)$ for $r=0,1,2$, and ST, VB to their starts. For each $Y$:
read the VB bit; copy $2^f$ records from each of E0, E1, E2 into BUF, which then holds
$\tilde x$ on $[J'2^f,(J'+3)2^f)$; read the $\Sigma$-record state. If valid, run $K_J$ with
$D_J$, step ST back $\Sigma$ and write the new state. Write the $2^f$ outputs (zeros if invalid
or $k\ge m$) to OUT and the $\xi$-record auxiliary row to AX.
\item Copy ST to FIN.
\end{enumerate}

\emph{Invariant.} At the start of $(J,Y)$, head E$r$ is at row $(J'(J)+r,Y)$, ST at record $Y$,
VB at bit $Y$, and OUT, AX at position $(X,J,Y)$ of slab order. It holds at $Y=0$ by the seeks,
and each step advances E$r$ by one row, ST by $\Sigma$, VB by one, OUT by $2^f$ and AX by $\xi$.

\emph{Correctness.} By induction on $J$, state record $Y$ of group $X$ equals the linewise
state of line $(X,Y)$ after slabs $0..J-1$: only step $(J,Y)$ touches record $Y$, and the slabs
of each line are visited in increasing $J$. By Claims~\ref{cl:window} and \ref{cl:splice}, BUF
contains the exact values the linewise kernel reads, so outputs, states and auxiliary rows
agree. Invalid lines keep a zero state and get zero outputs, as the linewise rule writes. No
padded box entry $j\in[\pi,n)$ is read (Claim~\ref{cl:splice}), and every output position is
written once, so the box keeps its full size and slab order.

\emph{Movement per valid group.} IN: $\le6G+n_eB$ (entering the group, copy, return, four
cross-group seeks, the build reads). CP: $\le3G+n_eB$. Each E$r$: $n_eB$ writes, a return
after the build and one at the next group, and in the main loop $B$ per slab plus a re-seek of
at most $B$ (head ends at $J'+r+1$, needs $J'(J+1)+r\in[J'+r,J'+r+2]$), so $\le3n_eB+2n_sB$.
OUT: $G$, plus one final return of $V$ in total. ST: $4|Y|\Sigma$ per slab and $4|Y|\Sigma$ at
the ends. VB: $2|Y|$ per slab. AX: $|Y|\xi$ per slab plus a final return. Summing with
$G=n_sB$:
\[
 |X|B(16n_s+11n_e)+V+\tfrac{4\Sigma+2\xi+2}{2^f}V+4|X||Y|\Sigma ,
\]
which is the bound in (b). Invalid groups cost less.

\subsection{Backward schedule}\label{sec:bwd}
Input is the forward output Z (slab order), AX and FIN. For each valid $X$: copy its $|Y|\Sigma$
final states from FIN to ST; for $J$ from $\lfloor(m-1)/2^f\rfloor$ down to $0$, seek Z and the
output tape to slab $(X,J)$ and AX to row $(X,J,0)$; for each $Y$ read the slab row into BUF and
its auxiliary row, read the state, run the backward kernel over $u$ \emph{decreasing} inside
BUF, write the new state, and write the $2^f$ outputs in increasing $u$. Lines are independent,
so visiting $Y$ in increasing order inside a decreasing-$J$ sweep is legitimate; only the order
along each line matters. Invariant and induction are as in the forward case with $J$
descending, and AX row $(X,J,Y)$ is read at the position where it was written.

Movement: Z and the output each read or write $B$ and seek back $2B$ per slab, plus entering
the group ($\le2G$): $\le5G$ each, and one final return of the output ($V$). AX:
$3|Y|\xi$ per slab. ST: $\le4|Y|\Sigma$ per slab plus loading. VB: $2|Y|$ per slab. This gives
(c). If a backward-type kernel needs neighbouring slabs (the second $M^{-1}$ after the backward
solve), use the three-tape window of Section~\ref{sec:fwd} with $J$ descending: $J'$ is
non-increasing with steps at most $2$, and the forward bound applies.

\subsection{Per-slab data and setup}
$D_J$ consists of $J'(J)$, the centres $c(k)$ and fractional offsets for $k$ in slab $J$, the
block boundaries and $q$-counters, the list of wraps in slab $J$, and (for Woodbury) the
elimination factors of the wraps in slab $J$. None depends on $X$ or $Y$. They are produced once
per axis by incrementing the $q$-counters along $J$, $O(2^fp)$ bit operations per slab, and stored
on DJ in $J$ order; DJ is swept once per X-group. Since it is read once per group and not per
line, the total is $|X|\cdot|D|$, with $|D|$ the size of all $D_J$ (bounded in
Section~\ref{sec:instances}). This replaces the notes' per-chunk odometer test, which is no
longer needed: validity is one VB bit per slab row. A block of the line algorithm that
straddles a slab boundary is recomputed from BUF by both slabs; with blocks of length at most
$W\le2^f/4$, at most two partial blocks occur per slab row, so kernel work is at most
$(1+2W/2^f)\le3/2$ times linewise.

\subsection{Instances and the conditions (W1)--(W3)}\label{sec:instances}
\begin{itemize}
\item $\mathsf S'$, $A$, $B_0$, $E$: windowed, no state ($\Sigma=\xi=0$). $W=3m+3$ for
$\mathsf S'$ (notes), the halo of $E$.
\item $M^{-1}$ via PR~\#5's chirped correlation: each segment block has length at most
$\lambda=\lfloor1/\theta\rfloor+1$ and reads a halo $w$; take $W\ge\lambda+2w$ so every segment
meeting slab $J$ lies in BUF and is recomputed there. No state. The segment crossing the period
boundary (``stitched from the two line ends'' in \texttt{segmented-inverse.tex}) is read through
$\tilde x$, i.e.\ through the halo slabs $v=-1$ and $v\ge n_\pi$; no separate stitching.
\item Woodbury. $N$ is the same matrix on every line of axis $i$, so $C=I+V^{\mathsf T}M^{-1}U$
and all its elimination factors are line-independent and live in $D_J$: $O(w^2)$ words of
precision $P$ per wrap, $|D|=O((\theta n+1)w^2P/b)$ records per axis, read once per group, a
fraction $O(\theta w^2P/b)$ of $V$ even when $|Y|=1$. Per line only right-hand sides are
carried. Forward sweep: the state is the running reduced right-hand side of the current block
row ($2w$ words) and the cyclic-border accumulator for the last block row ($2w$ words), so
$\Sigma=4wP/b$; the auxiliary row stores the reduced right-hand side of each wrap in the slab,
$\xi=(\lceil\theta2^f\rceil+1)\,2wP/b$. After slab $n_s-1$ every line of the group solves its
last $2w\times2w$ border system (factor in $D$) for $x_{\rm last}$, kept in FIN. Backward sweep:
state $(x_{k+1},x_{\rm last})$, $4w$ words; it reads the auxiliary rows in reverse.
\item Conditions. With the recentred sub-blocks $P\le2Q+O(\log s)=O(b)$. Then (W3) needs
$2^f\ge C w$ and $\theta w\le c$; with $\alpha^2\approx Q/(48d)$, $w=O(\sqrt d)$ and
$\theta\approx1/(4dL^y)$, so $\theta w\to0$, and (W1) with $W\ge\lambda+2w=\Theta(dL^y)$ gives
$2^f=\Theta(dL^y)$, $f=O(\log L)$. (W2) holds since $s_i\ge\tfrac45 t_i\ge r/4\gg2^f$.
\end{itemize}

\subsection{Total}
Per axis: exposure and restoration $12fVb$; each forward-type pass at most $28Vb$ and each
backward-type pass at most $11Vb$ up to the terms $11|X|(n_e-n_s)B=O(V2^f/n)$,
$O(V(\Sigma+\xi)/2^f)$ and $O(|X||Y|\Sigma)$, all $O(V)$ by (W3) and $o(V)$ as $n_s\to\infty$;
the setup terms above are $o(Vb)$. Over $d$ axes this is $O(Vb\,d(\log L+P))$, as the notes
claim. \qed

\section{Changes to the notes' paragraph}
\begin{enumerate}
\item The centred window $J'-1..J'+1$ with $J'=\lfloor J/\sigma\rfloor$ misses inputs; anchor
$J'$ at the lowest input read and use $J',J'+1,J'+2$ (Claim~\ref{cl:window}), with
$2^f\ge4W+4$. $J'$ moves by $0$, $1$ or $2$ per slab (2 occurs for $B_0$).
\item The wrap splice needs two heads advancing in lockstep over $Y$ (Claim~\ref{cl:splice});
copying ``the boundary slabs'' with one head costs $\Theta(|Y|)$ per record. The tail splice
pairs slab $n_\pi$ with slab~$0$, and the halo slabs pair consecutive slabs shifted by
$\pi\bmod2^f$. (W2) $\pi\ge3\cdot2^f$ keeps every virtual slab on at most two physical slabs.
\item Per-line state for all $|Y|$ interleaved lines is an array in $Y$ order on its own tape,
swept once per slab; auxiliary data for the backward sweep is an array in $(X,J,Y)$ order. Both
must fit in a slab row, which is (W3), and this is what lower-bounds $2^f$ besides the window.
\item The Woodbury factors are line-independent; only right-hand sides are per line. Storing
factors per line (as ``per-line state'' suggests) would cost $\Theta(w^2)$ per wrap per line and
violate (W3) when $\theta w^2P/b$ is not small.
\item Per-chunk validity by an odometer over $Y$ is replaced by a precomputed bitmap read in
lockstep; the hypothesis ``chunk volume $\gg\ell^2$'' is not needed.
\end{enumerate}

\section{Machine check}
\texttt{scripts/check\_fine\_field.py} represents every array as a tape with one head and counts every head
step. It runs exposure, a forward pass (windowed periodic map with $k$-dependent weights, a
two-word recurrence state, sparse wrap records to AX), a backward pass seeded from FIN that
reads AX in reverse, and restoration, for twelve shapes: $s\in\{13,29,53,61,103,107,127\}$
(prime, so $s\bmod2^f\ne0$), $2^f\in\{4,8,16\}$, $|Y|\in\{1,\dots,32\}$ with invalid $Y$
coordinates, $|X|\in\{2,3,4\}$ with invalid groups, and expansion, compression and $s\to s$
maps. All outputs equal the linewise maps; every case satisfies the bounds (a)--(c). Measured
forward $12.8$--$25.8\,V$, backward $3.8$--$9.6\,V$, exposure exactly $6V$ per bit; on longer
lines ($n=256$--$1024$, $2^f=16$--$64$) forward $18.4$--$21.3\,V$ and backward $9.0$--$11.3\,V$,
independent of $|Y|\in\{1,16,64\}$.
\end{document}
